\documentclass[11pt,a4paper]{article}

% --- Codificación e idioma ---
\usepackage[utf8]{inputenc}
\usepackage[T1]{fontenc}
\usepackage[spanish,es-tabla]{babel}
\usepackage{lmodern}

% --- Geometría y formato general ---
\usepackage[margin=2.5cm]{geometry}
\usepackage{microtype}
\usepackage{parskip}
\usepackage{enumitem}
\setlist{topsep=2pt,itemsep=2pt,parsep=0pt}

% --- Tablas y diagramas ---
\usepackage{array}
\usepackage{booktabs}
\usepackage{tikz}
\usetikzlibrary{positioning,arrows.meta}

% --- Cabeceras y numeración ---
\usepackage{fancyhdr}
\pagestyle{fancy}
\fancyhf{}
\fancyhead[L]{\small Bases de Datos II}
\fancyhead[R]{\small Resumen Teórico --- APIs / JDBC}
\fancyfoot[C]{\thepage}
\renewcommand{\headrulewidth}{0.4pt}

% --- Color y código ---
\usepackage{xcolor}
\definecolor{codebg}{RGB}{248,248,248}
\definecolor{codeframe}{RGB}{200,200,200}
\definecolor{codekw}{RGB}{0,0,180}
\definecolor{codestr}{RGB}{163,21,21}
\definecolor{codecmt}{RGB}{0,128,0}
\definecolor{notebg}{RGB}{255,250,230}
\definecolor{noteframe}{RGB}{220,180,80}

\usepackage{listings}

\lstdefinestyle{javastyle}{
    language=Java,
    basicstyle=\ttfamily\footnotesize,
    keywordstyle=\color{codekw}\bfseries,
    stringstyle=\color{codestr},
    commentstyle=\color{codecmt}\itshape,
    backgroundcolor=\color{codebg},
    frame=single,
    rulecolor=\color{codeframe},
    breaklines=true,
    breakatwhitespace=true,
    columns=fullflexible,
    keepspaces=true,
    showstringspaces=false,
    captionpos=b,
    xleftmargin=0.4em,
    xrightmargin=0.4em,
    aboveskip=8pt,
    belowskip=8pt,
    tabsize=2
}

\lstset{style=javastyle}

% --- Cajas de nota ---
\usepackage[most]{tcolorbox}
\newtcolorbox{nota}{colback=notebg,colframe=noteframe,arc=2pt,
    left=6pt,right=6pt,top=4pt,bottom=4pt,boxrule=0.6pt}

% --- Hipervínculos ---
\usepackage[hidelinks,bookmarks=true]{hyperref}

\title{Resumen Teórico --- APIs para Acceso a Bases de Datos\\[4pt]
       \large Bases de Datos II --- Práctico N.\textsuperscript{o} 4}
\author{Tomás Rodeghiero}
\date{2026}

\begin{document}

% =====================================================
% Carátula
% =====================================================
\begin{titlepage}
    \centering
    \vspace*{1.5cm}

    {\large \textbf{Universidad Nacional de Río Cuarto}}\\[6pt]
    {\normalsize Facultad de Ciencias Exactas, Físico-Químicas y Naturales}\\[4pt]
    {\normalsize Departamento de Computación}\\[4pt]
    {\normalsize Licenciatura en Ciencias de la Computación}\\[3cm]

    {\Large \textbf{Bases de Datos II}}\\[1.2cm]

    {\Large \textbf{Resumen Teórico}}\\[6pt]
    {\large APIs para Acceso a Bases de Datos --- JDBC\\[2pt]
    Práctico N.\textsuperscript{o} 4}\\[3cm]

    \begin{tabular}{rl}
        \textbf{Alumno:} & Tomás Rodeghiero \\[4pt]
        \textbf{Año:}    & 2026 \\
    \end{tabular}

    \vfill
    {\small Material elaborado a partir del teórico
    \emph{Teorico\_5\_APIs\_Accesos\_Base\_de\_Datos}.}
\end{titlepage}

\tableofcontents
\newpage

% =====================================================
\section{Introducción}
% =====================================================

Hasta este punto del cuatrimestre, todo lo que ejecutábamos sobre la
base de datos lo hacíamos con un cliente como \texttt{psql},
\texttt{mysql} o \texttt{sqlplus}: un humano escribiendo SQL.
El Práctico 4 cambia el ángulo: \emph{programación cliente} de la
base. Las aplicaciones reales no las escribe un humano comando a
comando: las escribe un programa, casi siempre en un lenguaje de
propósito general (Java, Python, C\#, etc.) usando una \textbf{API de
acceso a la base}.

El teórico se centra en la API de Java, \textbf{JDBC} (\emph{Java
DataBase Connectivity}), porque es la que se usa en el práctico, pero
las ideas son generales: hace falta un \emph{driver} específico del
motor, una \emph{conexión}, un objeto que represente la
\emph{sentencia} y una manera de \emph{recorrer los resultados}.

% =====================================================
\section{Arquitecturas de aplicaciones con acceso a base}
% =====================================================

Hay dos arquitecturas básicas:

\begin{description}[leftmargin=2.4cm,style=nextline]
    \item[\textbf{Dos capas}] La aplicación cliente se conecta
    \emph{directamente} al servidor de base de datos. Es lo más
    simple y lo que se usa en los ejercicios del práctico: una
    aplicación de consola Java con JDBC contra MySQL/PostgreSQL/Oracle.

    \item[\textbf{Tres capas}] La aplicación cliente se conecta a un
    \emph{servidor de aplicaciones} y éste, a su vez, dialoga con la
    base. Se usa en sistemas más grandes (web, servicios) por motivos
    de seguridad, escalabilidad y reutilización de lógica.
\end{description}

% =====================================================
\section{¿Qué es JDBC?}
% =====================================================

\textbf{JDBC} (\emph{Java DataBase Connectivity}) es una API Java que
permite ejecutar operaciones sobre bases de datos relacionales,
\emph{independientemente del motor concreto}, usando SQL.

Características principales:

\begin{itemize}
    \item Interacción con la base mediante SQL.
    \item 100\% Java (no requiere bindings nativos).
    \item Simple de usar: una decena de clases e interfaces alcanza.
    \item Alto rendimiento.
    \item Funciona contra cualquier base que provea un \emph{driver
    JDBC}: MySQL, PostgreSQL, Oracle, Firebird, etc.
\end{itemize}

% -----------------------------------------------------
\subsection{Arquitectura JDBC}
% -----------------------------------------------------

\begin{center}
\begin{tikzpicture}[node distance=0.8cm and 1.2cm,
    box/.style={rectangle,draw,minimum width=3.6cm,minimum height=0.8cm,
                font=\small},
    db/.style={cylinder,shape border rotate=90,aspect=0.4,draw,
               minimum width=2.4cm,minimum height=1.6cm,font=\small,
               cylinder uses custom fill,cylinder body fill=white},
    arr/.style={-{Latex},thick}]
    \node[box] (app) {Aplicación Java};
    \node[box,below=of app]  (api)  {API JDBC};
    \node[box,below=of api]  (drv)  {Driver de la base particular};
    \node[db,below=1.0cm of drv]    (db) {Base de Datos};
    \draw[arr] (app) -- (api);
    \draw[arr] (api) -- (drv);
    \draw[arr] (drv) -- (db);
\end{tikzpicture}
\end{center}

La aplicación Java conoce \emph{sólo} la API JDBC; el driver concreto
traduce esa API al protocolo nativo del motor. Cambiar de motor
cambia el driver y el URL, pero no las clases que usa la aplicación.

% =====================================================
\section{Componentes principales de la API}
% =====================================================

JDBC se construye alrededor de unos pocos componentes:

\begin{center}
\begin{tabular}{@{}p{4.4cm} p{9.5cm}@{}}
\toprule
\textbf{Componente} & \textbf{Rol} \\
\midrule
\texttt{DriverManager} (clase) &
Carga el driver de la base elegida y crea conexiones. \\
\texttt{Connection} (interfaz) &
Representa una sesión activa contra la base. \\
\texttt{Statement} y subclases (\texttt{PreparedStatement},
\texttt{CallableStatement}) &
Representan una sentencia SQL para ejecutar. \\
\texttt{ResultSet} (interfaz) &
Representa el conjunto de filas devueltas por una consulta. \\
\texttt{DatabaseMetaData} (interfaz) &
Provee información de metadatos: tablas, columnas, claves, etc. \\
\bottomrule
\end{tabular}
\end{center}

% =====================================================
\section{URL JDBC}
% =====================================================

Cada conexión JDBC se identifica con una URL del tipo:

\begin{lstlisting}
jdbc:subprotocolo:fuente
\end{lstlisting}

\begin{itemize}
    \item Cada driver tiene su propio \emph{subprotocolo}.
    \item Cada subprotocolo tiene su propia sintaxis para la
    \emph{fuente}.
\end{itemize}

\paragraph{Ejemplos.}

\begin{lstlisting}
jdbc:postgresql://host[:port]/database
jdbc:postgresql://localhost:5432/postgres

jdbc:mysql://host[:port]/database
jdbc:mysql://localhost:3306/practico1a_mysql

jdbc:oracle:thin:@//host:port/SID
jdbc:oracle:thin:@//localhost:1521/XE
\end{lstlisting}

% =====================================================
\section{Carga del driver y conexión}
% =====================================================

\subsection*{Clase \texttt{DriverManager}}

Provee el método estático:

\begin{lstlisting}
DriverManager.getConnection(String url, String user, String password);
\end{lstlisting}

Devuelve un objeto que implementa \texttt{Connection} y representa
una sesión con la base. El \emph{driver} debe estar previamente
registrado; con \texttt{Class.forName(\dots)} se fuerza la carga
explícita:

\begin{lstlisting}
// Ejemplo: conexion a MySQL.
String driver   = "com.mysql.cj.jdbc.Driver";
String url      = "jdbc:mysql://localhost/prueba";
String username = "root";
String password = "root";

// 1) Carga del driver (si todavia no esta cargado).
Class.forName(driver);

// 2) Apertura de la conexion.
Connection connection =
    DriverManager.getConnection(url, username, password);
\end{lstlisting}

\begin{nota}
\textbf{¿Es obligatorio \texttt{Class.forName}?} Desde JDBC 4 los
drivers se autorregistran vía \texttt{ServiceLoader}, así que en
muchos casos basta con tener el JAR en el \emph{classpath}. Aun así,
el teórico y la mayoría de los ejemplos prácticos siguen usando
\texttt{Class.forName} porque es la forma explícita y portable.
\end{nota}

% =====================================================
\section{Interfaz \texttt{Connection}}
% =====================================================

Una \texttt{Connection} representa una sesión con la base. Permite:

\begin{itemize}
    \item ejecutar instrucciones SQL y obtener resultados;
    \item obtener información de metadata (estructura) mediante
    \texttt{getMetaData()};
    \item manejar transacciones con \texttt{commit()} y
    \texttt{rollback()}.
\end{itemize}

\subsection*{Métodos para crear sentencias}

\begin{lstlisting}
Statement         createStatement();
PreparedStatement prepareStatement(String sql);
CallableStatement prepareCall(String sql);
\end{lstlisting}

\paragraph{¿Por qué tres tipos?} Por optimización y por casos de
uso: el \texttt{Statement} es el más genérico, el
\texttt{PreparedStatement} permite parámetros y precompilación, y el
\texttt{CallableStatement} se usa para invocar procedimientos
almacenados.

\subsection*{Transacciones: \texttt{setAutoCommit}}

\begin{itemize}
    \item Si \texttt{autoCommit} es \emph{verdadero} (valor por
    defecto), \emph{cada} sentencia ejecutada es una transacción
    independiente.
    \item Si se pasa a \emph{falso}, hay que delimitar la transacción
    explícitamente con \texttt{commit()} o \texttt{rollback()}.
\end{itemize}

\begin{lstlisting}
connection.setAutoCommit(false);
try {
    // varias sentencias...
    connection.commit();
} catch (SQLException e) {
    connection.rollback();
    throw e;
}
\end{lstlisting}

% =====================================================
\section{Tipos de \texttt{Statement}}
% =====================================================

\subsection*{\texttt{Statement}}

Representa una sentencia SQL \emph{estática}: no admite parámetros.
Útil para sentencias armadas dinámicamente que no reciben datos del
usuario, como las de exploración de metadatos.

\subsection*{\texttt{PreparedStatement}}

Hereda de \texttt{Statement} y representa una sentencia
\emph{precompilada} con parámetros (placeholders \texttt{?}).
Ventajas:

\begin{itemize}
    \item \textbf{Rendimiento}: el motor analiza y planifica una sola
    vez; las ejecuciones siguientes reutilizan el plan.
    \item \textbf{Seguridad}: como los parámetros viajan por separado
    de la sentencia, evita \emph{SQL injection}.
    \item \textbf{Conversión de tipos}: el driver se encarga de
    convertir \texttt{int}, \texttt{String}, \texttt{Date}, etc. al
    tipo SQL correspondiente.
\end{itemize}

\subsection*{\texttt{CallableStatement}}

Hereda de \texttt{Statement}; su objetivo es la \emph{ejecución de
procedimientos almacenados}. La sintaxis estándar (\emph{escape
call}) es:

\begin{lstlisting}
{ ? = call funcion(?, ?) }   // funcion con valor de retorno
{ call procedimiento(?, ?) } // procedimiento sin retorno
\end{lstlisting}

Para parámetros \texttt{OUT} o \texttt{INOUT} hay que registrarlos
con \texttt{registerOutParameter(int pos, int sqlType)}.

% =====================================================
\section{Parámetros en la creación de \texttt{Statement}}
% =====================================================

\texttt{Connection} ofrece versiones de \texttt{createStatement} y
\texttt{prepareStatement} que toman parámetros adicionales para
controlar el \texttt{ResultSet} resultante.

\subsection*{Tipo de \texttt{ResultSet}}

\begin{itemize}
    \item \texttt{ResultSet.TYPE\_FORWARD\_ONLY} (default): el cursor
    sólo avanza hacia adelante.
    \item \texttt{ResultSet.TYPE\_SCROLL\_INSENSITIVE}: el cursor se
    mueve en ambos sentidos; \emph{no} ve actualizaciones posteriores
    a la consulta.
    \item \texttt{ResultSet.TYPE\_SCROLL\_SENSITIVE}: ambos sentidos
    y refleja las actualizaciones cuando se producen.
\end{itemize}

\subsection*{Concurrencia / actualizabilidad}

\begin{itemize}
    \item \texttt{ResultSet.CONCUR\_READ\_ONLY} (default): sólo
    lectura.
    \item \texttt{ResultSet.CONCUR\_UPDATABLE}: el resultado es
    actualizable desde el propio \texttt{ResultSet}
    (\texttt{updateXxx}, \texttt{insertRow}, \texttt{deleteRow}).
\end{itemize}

% =====================================================
\section{Ejecución de sentencias}
% =====================================================

Métodos principales de \texttt{Statement}:

\begin{itemize}
    \item \texttt{ResultSet executeQuery(String sql)}: para
    \texttt{SELECT}; devuelve un \texttt{ResultSet}.
    \item \texttt{int executeUpdate(String sql)}: para
    \texttt{INSERT}, \texttt{UPDATE} o \texttt{DELETE}; devuelve la
    cantidad de filas afectadas.
    \item \texttt{boolean execute(String sql)}: para sentencias que
    pueden devolver múltiples resultados; \texttt{true} si el primer
    resultado es \texttt{ResultSet}.
\end{itemize}

% =====================================================
\section{Interfaz \texttt{ResultSet}}
% =====================================================

Un \texttt{ResultSet} provee acceso a la tabla generada por la
ejecución de una consulta. Conceptualmente es un puntero (cursor)
que apunta a una fila a la vez.

\subsection*{Recorrido y cierre}

\begin{itemize}
    \item \texttt{boolean next()}: avanza el cursor a la próxima
    fila. La \emph{primera} llamada activa la primera fila. Devuelve
    \texttt{false} cuando se queda sin filas.
    \item \texttt{void close()}: libera el \texttt{ResultSet} (en la
    práctica, conviene usar \texttt{try-with-resources}).
\end{itemize}

\subsection*{Lectura de columnas}

\begin{lstlisting}
<Tipo> get<Tipo>(int columnIndex);    // por indice (empieza en 1)
<Tipo> get<Tipo>(String columnName);  // por nombre  (menos eficiente)
\end{lstlisting}

\paragraph{Ejemplos.}

\begin{lstlisting}
int    nro    = rs.getInt("nro_cliente");
String nombre = rs.getString("nombre");
java.math.BigDecimal monto = rs.getBigDecimal("monto");
java.sql.Date fecha = rs.getDate("fecha");
\end{lstlisting}

\subsection*{Modificación de columnas (cursor actualizable)}

\begin{lstlisting}
update<Tipo>(String columnLabel, <Tipo> valor);
int findColumn(String columnName);  // indice de la columna por nombre
\end{lstlisting}

\subsection*{Navegación (en \texttt{ResultSet} \emph{scrollable})}

\begin{itemize}
    \item \texttt{previous()}: una fila atrás.
    \item \texttt{first()}, \texttt{last()}: a la primera/última.
    \item \texttt{beforeFirst()}, \texttt{afterLast()}: antes de la
    primera o después de la última (sin posicionar en una fila
    válida).
    \item \texttt{relative(int n)}: salta \emph{n} filas relativas a
    la actual.
    \item \texttt{absolute(int n)}: salta a la fila \emph{n}
    absoluta.
\end{itemize}

\subsection*{Ejemplo completo}

\begin{lstlisting}
Statement statement = connection.createStatement();
String query = "SELECT * FROM persona";

// Envia el query a la base y almacena el resultado.
ResultSet resultSet = statement.executeQuery(query);

// Recorre los resultados.
while (resultSet.next()) {
    System.out.print(" DNI: "    + resultSet.getString("DNI"));
    System.out.print("; Nombre: " + resultSet.getString("nombre"));
    System.out.print("; Email: "  + resultSet.getString("email"));
    System.out.println();
}
\end{lstlisting}

% =====================================================
\section{Extracción de metadatos: \texttt{DatabaseMetaData}}
% =====================================================

Con \texttt{Connection.getMetaData()} se obtiene un objeto
\texttt{DatabaseMetaData} que expone la \emph{estructura} de la base.
Es la pieza central del Ejercicio 3 del práctico.

\subsection*{Métodos más usados}

\paragraph{\texttt{getTables}.}

\begin{lstlisting}
ResultSet getTables(String catalog, String schemaPattern,
                    String tableNamePattern, String[] types);
\end{lstlisting}

Devuelve un \texttt{ResultSet} con al menos las columnas:

\begin{itemize}
    \item \texttt{TABLE\_CAT} (catálogo, puede ser \texttt{null}).
    \item \texttt{TABLE\_SCHEM} (esquema, puede ser \texttt{null}).
    \item \texttt{TABLE\_NAME} (nombre).
    \item \texttt{TABLE\_TYPE}: \texttt{TABLE}, \texttt{VIEW},
    \texttt{SYSTEM TABLE}, \texttt{GLOBAL TEMPORARY}, etc.
    \item \texttt{REMARKS} (comentarios).
\end{itemize}

\paragraph{\texttt{getColumns}.}

\begin{lstlisting}
ResultSet getColumns(String catalog, String schemaPattern,
                     String tableNamePattern, String columnNamePattern);
\end{lstlisting}

Cada fila describe una columna; el \texttt{ResultSet} tiene al menos
20 columnas. Los códigos de tipo de dato son los definidos en
\texttt{java.sql.Types}.

\paragraph{\texttt{getProcedures}.}

\begin{lstlisting}
ResultSet getProcedures(String catalog, String schemaPattern,
                        String procedureNamePattern);
\end{lstlisting}

Devuelve los procedimientos almacenados disponibles. Las constantes
\texttt{procedureColumnIn}, \texttt{procedureColumnOut}, etc.\ del
mismo \texttt{DatabaseMetaData} clasifican los parámetros.

\paragraph{Otros métodos útiles para el Ejercicio 3.}

\begin{itemize}
    \item \texttt{getPrimaryKeys(catalog, schema, table)}: devuelve
    las columnas de la PK con \texttt{KEY\_SEQ}.
    \item \texttt{getIndexInfo(catalog, schema, table, unique,
    approximate)}: devuelve los índices; combinado con un filtro
    sobre los que coinciden con la PK, sirve para detectar
    \emph{claves únicas adicionales}.
    \item \texttt{getImportedKeys(catalog, schema, table)} y
    \texttt{getExportedKeys(\dots)}: claves foráneas \emph{salientes}
    o \emph{entrantes}.
\end{itemize}

\subsection*{Ejemplo: listar esquemas en Oracle}

\begin{lstlisting}
String driver   = "oracle.jdbc.driver.OracleDriver";
String url      = "jdbc:oracle:thin:@//localhost:1521/XE";
String username = "USUARIO1";
String password = "USUARIO1";

Class.forName(driver);
Connection connection = DriverManager.getConnection(url, username, password);

DatabaseMetaData metaData = connection.getMetaData();
ResultSet resultSetSchemas = metaData.getSchemas();

System.out.println("Esquemas de la Base de datos");
while (resultSetSchemas.next()) {
    System.out.println("  esquema: " + resultSetSchemas.getString(1));
}
\end{lstlisting}

\subsection*{Ejemplo: listar tablas en PostgreSQL}

\begin{lstlisting}
String driver   = "org.postgresql.Driver";
String url      = "jdbc:postgresql://localhost:5432";
String username = "postgres";
String password = "root";

Class.forName(driver);
Connection connection = DriverManager.getConnection(url, username, password);

String[] tipo = { "TABLE" };
DatabaseMetaData metaData = connection.getMetaData();
ResultSet rs = metaData.getTables("Proyecto", "comparador", null, tipo);

while (rs.next()) {
    System.out.println("  catalogo: " + rs.getString(1));
    System.out.println("  esquema: "  + rs.getString(2));
    System.out.println("  nombre: "   + rs.getString(3));
    System.out.println("  tipo: "     + rs.getString(4));
    System.out.println("  comentarios: " + rs.getString(5));
}
\end{lstlisting}

% =====================================================
\section{Acceso al catálogo (cuando \texttt{DatabaseMetaData} no alcanza)}
% =====================================================

Hay información que la API \texttt{DatabaseMetaData} no expone
(detalles muy específicos de un motor). En esos casos se va
\emph{directamente al catálogo}: en MySQL existe la base
\texttt{information\_schema} y \texttt{mysql}; en PostgreSQL hay dos
catálogos por base, ANSI y PostgreSQL; en Oracle se usan las vistas
\texttt{USER\_*}/\texttt{ALL\_*}/\texttt{DBA\_*}.

\paragraph{Ejemplo (MySQL).} Listar nombres y código de los
procedimientos del esquema \texttt{procedimientos}:

\begin{lstlisting}
SELECT specific_name AS nombre, routine_definition AS codigo
  FROM information_schema.ROUTINES R
 WHERE routine_schema = 'procedimientos';
\end{lstlisting}

% =====================================================
\section{Correspondencia entre tipos SQL y Java}
% =====================================================

JDBC define una correspondencia estándar entre los tipos SQL y los
tipos Java; los métodos \texttt{getXxx}/\texttt{setXxx} respetan esa
tabla. Los datos que cruzan la red usan internamente
\texttt{java.sql.Types} como contrato.

\begin{center}
\begin{tabular}{@{}l l@{}}
\toprule
\textbf{Tipo SQL} & \textbf{Tipo Java} \\
\midrule
\texttt{CHAR}, \texttt{VARCHAR}, \texttt{LONGVARCHAR} &
\texttt{String} \\
\texttt{NUMERIC}, \texttt{DECIMAL} & \texttt{java.math.BigDecimal} \\
\texttt{BIT} & \texttt{boolean} \\
\texttt{TINYINT} & \texttt{byte} \\
\texttt{SMALLINT} & \texttt{short} \\
\texttt{INTEGER} & \texttt{int} \\
\texttt{BIGINT} & \texttt{long} \\
\texttt{REAL} & \texttt{float} \\
\texttt{FLOAT}, \texttt{DOUBLE} & \texttt{double} \\
\texttt{BINARY}, \texttt{VARBINARY}, \texttt{LONGVARBINARY} &
\texttt{byte[]} \\
\texttt{DATE} & \texttt{java.sql.Date} \\
\texttt{TIME} & \texttt{java.sql.Time} \\
\texttt{TIMESTAMP} & \texttt{java.sql.Timestamp} \\
\bottomrule
\end{tabular}
\end{center}

\begin{nota}
\textbf{\texttt{NUMERIC}/\texttt{DECIMAL} $\to$ \texttt{BigDecimal}.}
Por eso en el Ejercicio 2 del práctico la cotización máxima se lee
con \texttt{getBigDecimal(1)}: la columna SQL es
\texttt{NUMERIC(12,4)}.
\end{nota}

% =====================================================
\section{Persistencia de objetos en Java}
% =====================================================

Más allá de JDBC, la plataforma Java define APIs para
\emph{persistir objetos} en bases relacionales u objeto-relacionales:
\textbf{JDO} (\emph{Java Data Objects}), \textbf{JPA} (\emph{Java
Persistence API}), etc. Sus implementaciones (Hibernate, EclipseLink,
JPOX, OpenJPA\dots) terminan apoyándose en JDBC para hablar con el
motor.

\begin{center}
\begin{tikzpicture}[node distance=0.8cm and 1.2cm,
    box/.style={rectangle,draw,minimum width=4.2cm,minimum height=0.8cm,
                font=\small},
    db/.style={cylinder,shape border rotate=90,aspect=0.4,draw,
               minimum width=2.4cm,minimum height=1.6cm,font=\small,
               cylinder uses custom fill,cylinder body fill=white},
    arr/.style={-{Latex},thick}]
    \node[box] (app) {Aplicación Java};
    \node[box,below=of app]  (api)  {API (JDO, JPA, \dots)};
    \node[box,below=of api]  (impl) {Impl. de la API (JPOX, Hibernate)};
    \node[box,below=of impl] (jdbc) {Driver JDBC};
    \node[db,below=1.0cm of jdbc]   (db) {DB relacional};
    \draw[arr] (app)  -- node[right,font=\scriptsize]{utiliza}    (api);
    \draw[arr] (impl) -- node[right,font=\scriptsize]{implementa} (api);
    \draw[arr] (impl) -- node[right,font=\scriptsize]{utiliza}    (jdbc);
    \draw[arr] (jdbc) -- node[right,font=\scriptsize]{accede}     (db);
\end{tikzpicture}
\end{center}

% =====================================================
\section{Patrón típico con manejo de recursos}
% =====================================================

Los recursos JDBC son sensibles: si quedan abiertos pueden agotar el
\emph{pool} del motor. Desde Java 7, el patrón estándar es
\texttt{try-with-resources}, que cierra automáticamente cualquier
\texttt{AutoCloseable}.

\begin{lstlisting}
String sql = "SELECT id, nombre FROM persona WHERE id = ?";

try (Connection connection = DriverManager.getConnection(url, user, pass);
     PreparedStatement preparedStatement = connection.prepareStatement(sql)) {

    preparedStatement.setInt(1, 42);

    try (ResultSet resultSet = preparedStatement.executeQuery()) {
        while (resultSet.next()) {
            int    id     = resultSet.getInt("id");
            String nombre = resultSet.getString("nombre");
            // ...
        }
    }
}
\end{lstlisting}

% =====================================================
\section{Conexión con el Práctico 4}
% =====================================================

\begin{itemize}
    \item \textbf{Ejercicio 1}: usa \texttt{DriverManager},
    \texttt{Connection}, \texttt{PreparedStatement} (con
    \texttt{setInt}/\texttt{setString}/\texttt{executeUpdate}/
    \texttt{executeQuery}) y \texttt{ResultSet} con
    \texttt{next()}/\texttt{getXxx}. Es el caso ABM clásico.
    \item \textbf{Ejercicio 2}: agrega \texttt{CallableStatement} para
    invocar la función \texttt{fn\_actualizar\_cotizacion\_maxima} y
    leer un parámetro \texttt{OUT} con
    \texttt{registerOutParameter(\dots, Types.NUMERIC)} y
    \texttt{getBigDecimal(\dots)}. Suma un \emph{fallback} con
    \texttt{SELECT} para entornos PostgreSQL donde el escape
    \texttt{\{ ? = call \dots \}} no funciona.
    \item \textbf{Ejercicio 3}: usa \texttt{DatabaseMetaData} con
    \texttt{getTables}, \texttt{getColumns}, \texttt{getPrimaryKeys},
    \texttt{getIndexInfo} y \texttt{getImportedKeys}. Es
    independiente del motor (PostgreSQL u Oracle) cambiando sólo el
    \texttt{.properties}.
\end{itemize}

% =====================================================
\section{Ideas clave para repasar antes del parcial}
% =====================================================

\begin{enumerate}
    \item \textbf{JDBC es la API Java estándar de acceso a bases
    relacionales}; cada motor aporta su \emph{driver} concreto.
    \item \textbf{URL JDBC}: \texttt{jdbc:subprotocolo:fuente}.
    Conviene tenerla externalizada en un \texttt{.properties}.
    \item \textbf{Pasos básicos de conexión}: \texttt{Class.forName}
    + \texttt{DriverManager.getConnection(url, user, pass)}.
    \item \textbf{\texttt{Connection}} representa una sesión y es la
    fábrica de \texttt{Statement}, \texttt{PreparedStatement} y
    \texttt{CallableStatement}.
    \item \textbf{\texttt{Statement}} (estático), \emph{vs.}
    \textbf{\texttt{PreparedStatement}} (parametrizado, precompilado,
    seguro), \emph{vs.} \textbf{\texttt{CallableStatement}}
    (procedimientos almacenados).
    \item \textbf{Métodos de ejecución}: \texttt{executeQuery} (devuelve
    \texttt{ResultSet}), \texttt{executeUpdate} (devuelve
    \texttt{int}), \texttt{execute} (genérico).
    \item \textbf{\texttt{ResultSet}} es un cursor sobre filas;
    \texttt{next()} lo avanza, \texttt{getXxx} lee columnas (por
    índice o por nombre).
    \item \textbf{Tipos de \texttt{ResultSet}}: \texttt{FORWARD\_ONLY}
    (default), \texttt{SCROLL\_INSENSITIVE}, \texttt{SCROLL\_SENSITIVE}.
    Concurrencia: \texttt{CONCUR\_READ\_ONLY} o
    \texttt{CONCUR\_UPDATABLE}.
    \item \textbf{Transacciones}: \texttt{setAutoCommit(false)} +
    \texttt{commit()}/\texttt{rollback()}; con \texttt{autoCommit}
    \texttt{true} cada sentencia es su propia transacción.
    \item \textbf{Procedimientos almacenados}: \emph{escape call}
    \texttt{\{ ? = call f(?) \}}, parámetros \texttt{IN} con
    \texttt{setXxx}, \texttt{OUT} con \texttt{registerOutParameter}.
    \item \textbf{\texttt{DatabaseMetaData}} expone tablas, columnas,
    PK, FK, índices, procedimientos, esquemas. Es la herramienta del
    Ejercicio 3.
    \item \textbf{Correspondencia de tipos}: \texttt{NUMERIC} $\to$
    \texttt{BigDecimal}, \texttt{INTEGER} $\to$ \texttt{int},
    \texttt{VARCHAR} $\to$ \texttt{String}, \texttt{DATE} $\to$
    \texttt{java.sql.Date}, \texttt{TIMESTAMP} $\to$
    \texttt{java.sql.Timestamp}.
    \item \textbf{\texttt{try-with-resources}} para
    \texttt{Connection}, \texttt{Statement} y \texttt{ResultSet}: nunca
    se debe olvidar el \texttt{close()}.
    \item \textbf{Si \texttt{DatabaseMetaData} no alcanza}, se va al
    catálogo del motor (\texttt{information\_schema} en
    MySQL/PostgreSQL, vistas \texttt{USER\_*} en Oracle).
    \item \textbf{Por encima de JDBC} viven las APIs de
    \emph{persistencia de objetos} (JDO, JPA), pero su implementación
    sigue dialogando con la base por JDBC.
\end{enumerate}

% =====================================================
\section*{Resumen final}
\addcontentsline{toc}{section}{Resumen final}
% =====================================================

JDBC es la pieza que conecta el código de aplicación con la base. Su
arquitectura es minimalista pero potente: un \emph{driver} por motor,
una \emph{conexión} por sesión, distintos tipos de \emph{statement}
según el caso de uso, y un \emph{ResultSet} para recorrer filas.
\texttt{DatabaseMetaData} agrega una capa de introspección que es
útil cuando hace falta conocer la estructura sin recurrir al
catálogo del motor.

Si quedan firmes los cuatro bloques --- \emph{conexión}, \emph{tipo de
statement}, \emph{recorrido de \texttt{ResultSet}} y \emph{cierre con
\texttt{try-with-resources}} ---, el Práctico 4 se traduce
directamente: cada uno de los tres ejercicios es una variación sobre
un mismo esqueleto, intercambiando \texttt{PreparedStatement},
\texttt{CallableStatement} y \texttt{DatabaseMetaData} según lo que
pida la consigna.

\end{document}
